\documentclass[12pt,a4paper]{article}
\setcounter{secnumdepth}{3}
\setcounter{tocdepth}{3}

\usepackage{cmap}
\usepackage{amsmath}
\usepackage{amsfonts}
\usepackage{mathtext}
\usepackage{titlesec}
\usepackage{float}
\usepackage{subcaption}
\usepackage{amssymb}
\usepackage{tikz}
\usetikzlibrary{positioning, arrows.meta, shapes.geometric}
\usepackage[utf8]{inputenc}
\usepackage[T2A]{fontenc}
\usepackage{wrapfig}
\usepackage[english, russian]{babel}
\usepackage[left=2cm,right=2cm,top=2cm,bottom=2cm]{geometry}
\usepackage{indentfirst}
\usepackage{listings}
\usepackage{xcolor}

\lstset{
    basicstyle=\ttfamily\small,
    keywordstyle=\color{blue},
    commentstyle=\color{green!60!black},
    stringstyle=\color{red},
    numbers=left,
    numberstyle=\tiny\color{gray},
    stepnumber=1,
    numbersep=5pt,
    backgroundcolor=\color{gray!10},
    showspaces=false,
    showstringspaces=false,
    showtabs=false,
    frame=single,
    tabsize=2,
    captionpos=b,
    breaklines=true,
    breakatwhitespace=false,
    title=\lstname,
    escapeinside={\%*}{*)},
    morekeywords={*,...},
    language=Python
}

\DeclareSymbolFont{T2Aletters}{T2A}{cmr}{m}{it}

%%% Работа с картинками
\usepackage{graphicx}  % Для вставки рисунков
\graphicspath{{imgs/}}  % папки с картинками

\usepackage{caption}
\captionsetup{labelsep=period, labelfont=bf}

\titleformat{\section}[block]
{\bfseries}{\thesection.}{0 cm}{}

\titleformat{\subsection}[block]
{\bfseries}{\thesubsection.}{0 cm}{}

\titleformat{\subsubsection}[block]
{\bfseries}{\thesubsubsection.}{0 cm}{}

\titlespacing{\section}{17pt}{14pt}{10pt}
\titlespacing{\subsection}{17pt}{14pt}{4pt}
\titlespacing{\subsubsection}{17pt}{14pt}{4pt}

\def \TITLE {Отчет о выполнении лабораторной работы №1}
\def \SUBTITLE {Измерение температуры пламени методом
обращения спектральных линий}
\def \AUTHOR {Выполнили студенты группы Б03-405\\Подшивалова Анастасия\\Савельева Александра\\Тимохин Даниил\\Трифонова Анна\\Собкина Ксения}
\def \DATE {16 марта 2026 г.}

\begin{document}

\begin{titlepage}
	\centering
	\vspace{5cm}
	{\scshape\large Московский физико-технический институт \\
	(НАЦИОНАЛЬНЫЙ ИССЛЕДОВАТЕЛЬСКИЙ УНИВЕРСИТЕТ)}
	
	\vspace{4cm}
	{\LARGE \TITLE}
	
	\vspace{1cm}
	{\Huge\bf \SUBTITLE }
	
	\vspace{1cm}
	\vfill
	
\begin{flushright}
	{\LARGE \AUTHOR}
\end{flushright}
	

	\vfill

	\DATE
\end{titlepage}

\newpage

\fontsize{12}{14}\selectfont

\section{ Аннотация}


\section{ Теоретическая справка}


\section{ Проведение эксперимента и обработка результатов}
Сначала воспользуемся методом обращения спектральных линий для нахождения при каких параметрах тока япкость лампы совпадает с яркостью пламени.

После для каждого параметра тока надем параметры тока, при которых вольфрамовая нить сливается с излучением лампы.

\begin{table}[h]
\centering
\caption{Измерения}
\begin{tabular}{|c|c|c|}
\hline
\textbf{Ток лампы, A} & \textbf{Ток пирометра, дел} & \textbf{Ток пирометра, A} \\\hline
$12.1$   & $67$ & $0.335$ \\\hline
$16.3$   & $66$ & $0.33$ \\\hline
$17.4$   & $60$ & $0.34$ \\\hline
$16.8$   & $66$ & $0.33$ \\\hline
$17.2$   & $54$ & $0.27$ \\\hline
\end{tabular}
\end{table}

Сначала вычислим средний ток через спираль пирометра. $I=0.335$ А.

Теперь, имея данные о параметрах самой устанвки из формулы(), так как температура не мнеяется, но меняется исследуемая длина волны, то можно вычислить температуру пламени

\begin{equation}
    T=\frac{T_b C_2}{C_2+T_b(\lambda_ж - \lambda_к)\ln{\varepsilon}}
\end{equation}



\section{ Обсуждение результатов и выводы}


\end{document}
